\input{preamble}

\begin{document}

\title{Historia do Brasil}
\maketitle

\phantomsection
\label{section-phantom}

\tableofcontents

\section{Guerra do Paraguai}
\label{section-guerra-do-paraguai}

A Guerra do Paraguai ocorreu entre 1864 e
1870. Uma maneira simples de entender seu
início é separar duas perguntas: quem começou
a crise regional e quem atacou diretamente o
Brasil primeiro.

Em 1864, o Brasil interveio militarmente no
Uruguai contra o governo dos Blancos. O
presidente paraguaio Francisco Solano López,
aliado aos Blancos, considerava a intervenção
brasileira uma ameaça ao equilíbrio de poder
na região do Prata.

O Brasil, porém, não atacou o Paraguai nesse
momento. Foi o Paraguai que iniciou as
hostilidades diretas contra o Brasil: em
novembro de 1864, o governo paraguaio
apreendeu o vapor brasileiro \emph{Marquês de
Olinda}; em dezembro, tropas paraguaias
invadiram Mato Grosso.

Portanto, a sequência inicial pode ser
resumida assim: o Brasil intervém no Uruguai;
López responde atacando o Brasil. Dizer que
López simplesmente queria conquistar o Brasil
é enganoso. Havia ambições paraguaias de
poder regional, mas a guerra surgiu de uma
crise mais ampla envolvendo Brasil, Paraguai,
Argentina e Uruguai, com disputas de
fronteiras, navegação dos rios e influência
política na Bacia do Prata.

Em 1865, López tentou enviar tropas através
do território argentino para intervir no
Uruguai. O governo argentino recusou a
passagem, e o Paraguai invadiu Corrientes.
Nesse mesmo período, a intervenção brasileira
ajudou a derrubar o governo blanco uruguaio.
Venancio Flores, dos Colorados e aliado do
Brasil, assumiu o poder.

Forma-se então a Tríplice Aliança: Brasil,
Argentina e o novo governo uruguaio contra o
Paraguai. Há aqui uma ironia importante: López
entrou no conflito em parte para defender o
governo blanco do Uruguai contra a intervenção
brasileira; depois que esse governo caiu, o
novo Uruguai passou a combater ao lado do
Brasil contra López.

Uma cronologia mínima do começo da guerra é:

\begin{enumerate}
\item Brasil intervém militarmente no Uruguai
em 1864.
\item Paraguai apreende o \emph{Marquês de
Olinda} em novembro de 1864.
\item Paraguai invade Mato Grosso em dezembro
de 1864.
\item Paraguai invade Corrientes em 1865.
\item Brasil, Argentina e Uruguai formam a
Tríplice Aliança em 1865.
\end{enumerate}

Assim, à pergunta ``quem atacou quem
primeiro?'', a resposta precisa é: o Brasil
atacou primeiro no contexto da crise regional,
mas atacando o Uruguai; entre Brasil e
Paraguai, foi o Paraguai que atacou o Brasil
primeiro.

\end{document}
